\documentclass[../main.tex]{subfiles}
\begin{document}
\section{MACs y Cifrado Autenticado}

Hasta ahora la criptografía vista resolvía \textbf{confidencialidad} (controlar la
diseminación de información). Esta unidad ataca el otro gran problema: la
\textbf{integridad}, es decir, cómo confiar en que la información no fue
modificada, aun cuando el origen es confiable. La primitiva básica para esto es
el \textbf{MAC} (Message Authentication Code), y luego veremos funciones de hash y
cómo combinar confidencialidad con integridad (cifrado autenticado).

\prof{En general el problema de integridad no es evitar la modificación, porque muchas veces no se puede, sino detectarla.}

\subsection{El nuevo escenario de ataque: integridad}

El cifrado por sí solo no protege contra \textbf{modificaciones} del texto cifrado.
El profesor lo motiva con un sistema de RR.HH. que guarda sueldos cifrados en una
base de datos, con un criptosistema de flujo CPA-Secure.

\begin{itemize}[leftmargin=1.4em]
  \item \textbf{No se puede extraer el valor:} aunque el atacante sepa que un sueldo
  es \$1.000.000, como el criptosistema es \emph{no determinista} (usa IV), cifrar
  ese valor da algo distinto cada vez. No puede identificar qué fila es ni aproximar
  el sueldo. La parte de confidencialidad es realmente segura.
  \item \textbf{Pero sí se puede manipular:} un administrador que conoce su legajo
  (2678) y el de su jefe (2890) puede copiar el texto cifrado del jefe sobre el suyo
  ($update$) y obtener un aumento de sueldo, \emph{sin conocer ningún valor}. El
  escenario CPA se mantiene (nunca extrajo información), pero igual rompió el
  sistema.
\end{itemize}

\paragraph{Maleabilidad.} Es la propiedad de que una modificación del texto cifrado
se propaga de forma sistemática/controlable al texto plano.

\begin{definicion}{Maleabilidad}
Cuando una modificación en el texto cifrado se puede traducir de forma sistemática
en alguna modificación en el texto plano. La más simple: cambio un bit del cifrado
y cambia un bit del plano.
\end{definicion}

\paragraph{Ataque quirúrgico sobre cifrado de flujo.} Si $\mathrm{Enc}_k(m)=G(k)\oplus m=c$,
entonces definiendo $c'=c\oplus x$ resulta $\mathrm{Dec}_k(c')=m\oplus x$. Conociendo
el propio sueldo en plano ($m=1345$) y su cifrado, el atacante hace:
$$c'' = c \oplus (1345) \oplus (100000)$$
y al descifrarse da $1345 \oplus 1345 \oplus 100000 = 100000$. Cambió el sueldo a un
valor arbitrario aunque la información estuviera cifrada. Aclaración: agregar el
legajo dentro del cifrado (\texttt{empleado\textbar\textbar sueldo}) frena la copia
de columnas, pero \textbf{no} frena este ataque de maleabilidad.

\paragraph{CCA: el modelo formal.} Para modelar esto se define el escenario de
\textbf{texto cifrado escogido} (Chosen Ciphertext Attack, $\mathrm{CCA}_{A,\Pi}$).
Es como CPA pero el adversario además tiene un \textbf{oráculo de descifrado}
$g(x)=d_k(x)$ (no solo el de cifrado $f(x)=e_k(x)$). Restricción para no hacerlo
trivial: al recibir el desafío $c=e_k(m_b)$ \emph{no} puede pedir $g(c)$.

\begin{destacado}
\textbf{Ningún criptosistema visto (ni de los no vistos) es CCA-Secure.} El cifrado
está pensado para confidencialidad (transferencia de información), no para
manipulación. Por eso la integridad requiere otra primitiva: el \textbf{MAC}.
\end{destacado}

\begin{examen}
Ejercicio de cátedra: \textbf{demostrar que el cifrado de flujo en general no es
CCA-Secure.} Ayuda: considerar $m_0=(0\ldots0)$ y $m_1=(1\ldots1)$ y ver qué
consultas se pueden hacer a $g(x)$ al recibir $c$. La idea es usar la maleabilidad
($\oplus$) para transformar $c$ en otro cifrado consultable.
\end{examen}

\subsection{Qué es un MAC}

\begin{definicion}{MAC (Message Authentication Code)}
Terna de algoritmos $(\mathrm{Gen}, \mathrm{Mac}, \mathrm{Vrfy})$:
\begin{itemize}[leftmargin=1.4em,nosep]
  \item $\mathrm{Gen}: () \to K$ \quad — generador de clave.
  \item $\mathrm{Mac}: K \times P \to T$ \quad — etiquetador (genera la etiqueta/tag).
  \item $\mathrm{Vrfy}: K \times P \times T \to \{0,1\}$ \quad — verificador.
\end{itemize}
\textbf{Propiedad básica (correctitud):} para todo $m$ y $k$ válidos,
$\mathrm{Vrfy}_k(m, \mathrm{Mac}_k(m)) = 1$.
\end{definicion}

Los conjuntos: $K$ = espacio de claves, $P$ = espacio de mensajes, $T$ = espacio de
etiquetas. A diferencia de los criptosistemas, hay etiquetado y verificación (no
cifrado/descifrado). La verificación es parte del diseño del MAC; en muchos casos
recalcula la etiqueta y compara (\texttt{Mac\_k(m) == t}), pero no siempre (ver
ejemplo del cifrado de la longitud).

\begin{center}
\begin{tikzpicture}[
  font=\footnotesize,
  box/.style={draw,rounded corners,minimum height=8mm,minimum width=13mm,align=center},
  >={Stealth[]}]
  % Emisor
  \node (m) at (0,0) {$m$};
  \node[box,fill=blue!8] (mac) at (0,-1.3) {$\mathrm{Mac}_k$};
  \node (t) at (0,-2.6) {$t$};
  \draw[->] (m) -- (mac);
  \draw[->] (mac) -- (t);
  \node[above=1pt of m] {\textbf{Emisor}};

  % Canal: envia (m,t)
  \node[draw,dashed,rounded corners,minimum width=22mm,minimum height=7mm] (canal) at (3.3,-1.3) {$(m,\,t)$};
  \draw[->] (m) to[out=0,in=150] (canal.west);
  \draw[->] (t) to[out=0,in=210] (canal.west);

  % Receptor
  \node[box,fill=blue!8] (vrfy) at (6.6,-1.3) {$\mathrm{Vrfy}_k$};
  \node (out) at (9,-1.3) {$\{0,1\}$};
  \draw[->] (canal.east) -- (vrfy);
  \draw[->] (vrfy) -- (out) node[midway,above]{$1$/$0$};
  \node[above=1pt of vrfy] {\textbf{Receptor}};

  % clave compartida
  \node[draw,fill=yellow!25,rounded corners] (k) at (3.3,1) {clave compartida $k$};
  \draw[->,gray] (k) to[out=200,in=90] (mac.north);
  \draw[->,gray] (k) to[out=340,in=90] (vrfy.north);
\end{tikzpicture}
\end{center}
\diagcaption{El emisor calcula $t=\mathrm{Mac}_k(m)$ y envía $(m,t)$; el receptor corre $\mathrm{Vrfy}_k(m,t)$ y obtiene $1$ (íntegro) o $0$ (alterado). La misma clave $k$ se comparte entre ambos.}

\paragraph{Integridad = detectar, no evitar.} En un canal inseguro alguien puede
pararse en el medio; no se puede impedir la modificación. Lo que da el MAC es la
capacidad de \textbf{detectarla} para luego decidir qué hacer (retransmitir, alertar,
abortar).

\subsection{Seguridad de un MAC: infalsificabilidad (Mac-forge)}

\begin{definicion}{Experimento Mac-forge$_{A,\Pi}$}
Dado nivel de seguridad $n$, adversario $A$ y MAC $\Pi(n)$:
\begin{enumerate}[leftmargin=1.6em,nosep]
  \item Se genera una clave $k \leftarrow K$.
  \item $A$ obtiene un oráculo etiquetador $f(x)=\mathrm{Mac}_k(x)$ (caja negra: nunca ve $k$).
  \item $A$ realiza todas las evaluaciones que quiera. Sea $Q$ el conjunto de mensajes evaluados.
  \item $A$ emite $(m, t)$ con $m \notin Q$.
\end{enumerate}
$\mathrm{Mac\text{-}Forge}_{A,\Pi} = 1$ si $\mathrm{Vrfy}_k(m,t)=1$ (falsificación válida).\\
Si $\Pr[\mathrm{Mac\text{-}Forge}_{A,\Pi}=1] \le \mathrm{neg}(n)$ entonces $\Pi$ es
\textbf{infalsificable} (existencialmente infalsificable bajo ataque de mensaje
elegido, CMA).
\end{definicion}

La restricción $m \notin Q$ evita el ataque trivial: si pedí el tag de ``Hola'' no
puedo emitir ``Hola'' con ese tag. Una falsificación es generar, \emph{sin conocer la
clave}, un par $(m,t)$ que pase la verificación.

\paragraph{Observaciones clave.}
\begin{itemize}[leftmargin=1.4em]
  \item El adversario puede obtener el MAC de cualquier mensaje que elija.
  \item Se considera \textbf{roto} el MAC si el adversario puede falsificar cualquier
  mensaje, \emph{independientemente de si tiene sentido o no}. (Históricamente el
  ataque de Wang 2004 a MD5 generaba falsificaciones ``sin sentido'', pero al
  acelerarse y poder restringir la salida a caracteres imprimibles, pasó a ser un
  ataque real. Lección: un ataque debilitado no se debe tomar a la ligera.)
  \item Demostrar que algo es infalsificable es muy difícil (hay que cubrirse de todo
  adversario); demostrar que es \textbf{falsificable} es más fácil: basta exhibir
  \emph{un} adversario (contraejemplo).
\end{itemize}

\subsection{MÉTODO para romper un MAC inseguro}

\begin{examen}
\textbf{Pregunta recurrente (2015, 2017): romper un MAC vía Mac-forge.}\\
Método general para mostrar que un MAC es \textbf{falsificable} (= inseguro):
\begin{enumerate}[leftmargin=1.6em]
  \item Ubicarse como adversario $A$ con acceso al oráculo $f(x)=\mathrm{Mac}_k(x)$.
  \item Hacer una o pocas evaluaciones inteligentes (típicamente $f(00\ldots0)$ para
  ``leer'' material de la clave/keystream, o el MAC de un prefijo).
  \item \textbf{Exhibir una falsificación existencial:} construir un par $(m,t)$
  con $m \notin Q$ (distinto de lo evaluado) tal que $\mathrm{Vrfy}_k(m,t)=1$, sin
  haber consultado el tag de ese $m$.
  \item Verificar que el par propuesto pasa la verificación.
\end{enumerate}
\end{examen}

\paragraph{Tres MACs de juguete (ejercicio de cátedra) y por qué fallan.}

\textbf{(1) $\mathrm{Mac}_k(m) = G(k) \oplus m$} (usar un cifrado de flujo como MAC).
La verificación: $\mathrm{Vrfy}_k(m,t) = (t \oplus m == G(k))$. Es \emph{falsificable}:
$G(k)$ es determinístico, así que si averiguo $G(k)$ puedo etiquetar cualquier cosa.
\begin{itemize}[leftmargin=1.4em,nosep]
  \item Evalúo $f(00\ldots0)$. Como $G(k)\oplus 0 = G(k)$, obtengo directamente $G(k)$.
  \item Emito $m=11\ldots1$ (distinto de $0\ldots0$) con tag $t = G(k)\oplus(11\ldots1)$.
  \item Verifica: $t \oplus m = G(k)\oplus 1\ldots1 \oplus 1\ldots1 = G(k)$. Falsificación válida.
\end{itemize}
\emph{Lección:} el problema es el determinismo que se puede extraer en un mensaje y
reusar en otro.

\textbf{(2) $\mathrm{Mac}_k(m) = k \oplus \mathrm{first\_k\_bits}(m)$} (solo usa los
primeros bits del mensaje). Doblemente falsificable:
\begin{itemize}[leftmargin=1.4em,nosep]
  \item Evalúo $f(00\ldots0)$: como los primeros bits son 0, el tag es $k$. ¡Recuperé la clave!
  Luego emito cualquier $m$ con tag $k \oplus \mathrm{first\_k\_bits}(m)$.
  \item Más fácil aún: como solo mira \emph{parte} del mensaje, tomo un $m$, calculo su
  MAC, y emito $m$ con un carácter extra al final (la parte no chequeada). La etiqueta
  sigue siendo válida.
\end{itemize}
\begin{destacado}
Un MAC que \textbf{no procesa la totalidad del mensaje} es trivialmente falsificable:
modifico el bit/byte que quedó afuera y reuso la misma etiqueta. Un MAC robusto
\textbf{siempre} es resultado de procesar todo el contenido del mensaje.
\end{destacado}

\textbf{(3) $\mathrm{Mac}_k(m) = \mathrm{Enc}_k(|m|)$} (cifrar la \emph{longitud} con un
cripto CPA-Secure). Como $\mathrm{Enc}$ es no determinista, dos mensajes de la misma
longitud dan tags distintos (parece robusto desde afuera). La verificación debe
\emph{descifrar} el tag y comparar con $|m|$ (no se puede recifrar). Falsificable
porque \textbf{la etiqueta no depende del contenido}, solo de la longitud:
\begin{itemize}[leftmargin=1.4em,nosep]
  \item Pido el MAC de un mensaje de longitud 2 (p.ej. $00$); obtengo un tag $t$ tal que
  $\mathrm{Dec}_k(t)=2$.
  \item Emito otro mensaje distinto de la misma longitud (p.ej. $11$) con el \emph{mismo} $t$.
  Pasa la verificación porque al descifrar $t$ vuelve a dar $2 = |11|$.
  \item Equivalente: pedir el etiquetado de dos mensajes de igual longitud e
  \textbf{intercambiar las etiquetas} (ambas validan).
\end{itemize}

\begin{examen}
\textbf{Falsificación por bloques (parcial 2015):} dado
$\mathrm{Mac}_k(m)=\langle F_k(m_1), F_k(F_k(m_2))\rangle$ y similares, el método es
\textbf{reusar/intercambiar bloques} entre mensajes para forjar el tag de otro
mensaje (los bloques se calculan de forma independiente).\\[2pt]
\textbf{MAC por bit de paridad (parcial 2017):} si el tag es la paridad del mensaje,
\textbf{cambiar dos bits} del mensaje conserva la paridad $\Rightarrow$ el tag sigue
válido $\Rightarrow$ falsificación existencial.
\end{examen}

\subsection{Construcción de MAC: CBC-MAC}

A partir de una primitiva de \textbf{cifrado de bloque} (sea $F$ una función
pseudoaleatoria) se construye un MAC infalsificable. Existe CBC-MAC con cualquier
primitiva: 3DES, IDEA, AES, etc.

\begin{definicion}{CBC-MAC (longitud de mensaje fija)}
$\mathrm{Gen}: k \leftarrow \{0,1\}^n$.\\
$\mathrm{Mac}$: sea $m = m_1\|m_2\|\cdots\|m_j$, con $t_0 = 00\ldots0$ (IV en cero, fijo):
$$t_i = F_k(t_{i-1} \oplus m_i), \qquad \mathrm{Mac}_k(m) = t_j.$$
$\mathrm{Vrfy}$: $\mathrm{Vrfy}_k(m,t)=1 \iff t = \mathrm{Mac}_k(m)$ (se recalcula; $F_k$ es
determinística).
\end{definicion}

Es CBC pero \emph{sin} interés en recuperar el mensaje: se tiran todos los resultados
intermedios y se conserva solo el cifrado del último bloque, que ``depende'' de todos
los bloques anteriores por el encadenamiento. El IV se pone en cero (no estamos
cifrando, no necesita ser aleatorio).

\begin{center}
\begin{tikzpicture}[
  font=\footnotesize,
  arrow/.style={-{Stealth}},
  box/.style={draw,fill=blue!8,minimum height=8mm,minimum width=11mm},
  xor/.style={draw,circle,inner sep=1pt}]
  % bloques de mensaje
  \node (m1) at (0,0)   {$m_1$};
  \node (m2) at (2.6,0) {$m_2$};
  \node (m3) at (5.2,0) {$\cdots$};
  \node (ml) at (7.8,0) {$m_\ell$};
  % xors
  \node[xor] (x1) at (0,-1)   {$\oplus$};
  \node[xor] (x2) at (2.6,-1) {$\oplus$};
  \node[xor] (x3) at (5.2,-1) {$\oplus$};
  \node[xor] (xl) at (7.8,-1) {$\oplus$};
  % cajas F_k
  \node[box] (f1) at (0,-2.1)   {$F_k$};
  \node[box] (f2) at (2.6,-2.1) {$F_k$};
  \node[box] (f3) at (5.2,-2.1) {$F_k$};
  \node[box] (fl) at (7.8,-2.1) {$F_k$};
  % IV
  \node (iv) at (-2.2,-1) {$t_0=00\ldots0$};
  \draw[arrow] (iv) -- (x1);
  % entradas de mensaje
  \draw[arrow] (m1) -- (x1);
  \draw[arrow] (m2) -- (x2);
  \draw[arrow] (m3) -- (x3);
  \draw[arrow] (ml) -- (xl);
  % xor -> F_k
  \draw[arrow] (x1) -- (f1);
  \draw[arrow] (x2) -- (f2);
  \draw[arrow] (x3) -- (f3);
  \draw[arrow] (xl) -- (fl);
  % encadenamiento t_i hacia el siguiente xor
  \draw[arrow] (f1) -| (2.6,-1.5) node[midway,below]{$t_1$} -- (x2);
  \draw[arrow] (f2) -| (5.2,-1.5) node[midway,below]{$t_2$} -- (x3);
  \draw[arrow] (f3) -| (7.8,-1.5) node[midway,below]{$t_{\ell-1}$} -- (xl);
  % salida = tag
  \node[draw,fill=yellow!25,rounded corners] (tag) at (10.5,-2.1) {$t=\mathrm{Mac}_k(m)$};
  \draw[arrow] (fl) -- (tag) node[midway,above]{$t_\ell$};
\end{tikzpicture}
\end{center}
\diagcaption{CBC-MAC: con $t_0=00\ldots0$ (IV en cero), cada bloque hace $t_i=F_k(t_{i-1}\oplus m_i)$. El encadenamiento hace que el \emph{último} bloque $t_\ell$ dependa de todo el mensaje; se descartan los intermedios y la salida es $t_\ell$.}

\begin{examen}
\textbf{¿Bajo qué condiciones es seguro CBC-MAC? (parcial 2016)}\\
La construcción básica es infalsificable \textbf{SOLO si se permiten mensajes de una
misma longitud fija} (ni más grandes ni más chicos), con \textbf{IV constante/cero}.
Con longitud variable aparecen \textbf{ataques de extensión}.
\end{examen}

\paragraph{Ataque de extensión (longitud variable).} Idea de la falsificación:
\begin{itemize}[leftmargin=1.4em,nosep]
  \item Crear dos bloques aleatorios $A$ y $B$. Definir $m_1 = A$ (un bloque) y $m_2 = A\|B$.
  \item Calcular $t_1 = f(A)$ (= $F_k(A)$) y $t_2 = f(A\|B)$.
  \item Emitir el mensaje de \textbf{tres bloques} $A \,\|\, B \,\|\, (A \oplus t_1)$ con tag $t_2$.
\end{itemize}
\emph{Por qué valida:} el estado tras procesar los dos primeros bloques $A\|B$ es
$t_1$. El tercer bloque entra como $F_k(t_1 \oplus (A\oplus t_1)) = F_k(A) = t_1 = f(A)$...
es decir, abusando del encadenamiento se fuerza la entrada del bloque final a un valor
$A$ cuyo MAC ya se conoce. El truco general: construir mensajes por \textbf{prefijo}
para acceder a los \textbf{estados intermedios} que normalmente se descartan, y luego
anular el estado con $\oplus$ y forzar el valor deseado (muy parecido al ataque de
maleabilidad del cifrado de flujo).

\paragraph{Extensiones seguras de CBC-MAC (mensajes arbitrarios).} Tres formas (todas
seguras; ninguna es perfecta):
\begin{enumerate}[leftmargin=1.6em,nosep]
  \item \textbf{Clave dependiente de la longitud:} $k' = F_k(|m|)$ y usar $t_i = F_{k'}(t_{i-1}\oplus m_i)$.
  Así dos mensajes de distinta longitud usan claves distintas y no se aprenden los
  estados intermedios. (Requiere conocer $|m|$ a priori.)
  \item \textbf{Longitud como prefijo:} $m' = |m| \,\|\, m$, y $\mathrm{Mac}_k(m)=\mathrm{CBC\text{-}MAC}_k(m')$.
  Rompe que un mensaje sea prefijo de otro. (Requiere conocer $|m|$ a priori.)
  \item \textbf{Doble clave (recifrado del último estado):} $k_1, k_2$; calcular
  $t' = \mathrm{CBC\text{-}MAC}_{k_1}(m)$ y luego $t = F_{k_2}(t')$. No requiere conocer
  $|m|$ a priori, pero usa el doble de material de clave.
\end{enumerate}

\begin{ojo}
Poner la longitud como \textbf{sufijo} (procesar el mensaje y agregar $|m|$ al final)
\textbf{NO es seguro}: hay ataques de extensión parecidos.
\end{ojo}

\subsection{Funciones de hash}

\begin{definicion}{Función de hash}
Par de algoritmos: un selector que elige una función de una familia (información
\textbf{pública}, no es clave) y la función $h$ que convierte un mensaje arbitrario en
una etiqueta (``hash'' / digest / resumen) de tamaño \textbf{fijo y reducido}.
\textbf{No tiene clave} (diferencia fundamental con MAC y criptosistemas).
\end{definicion}

Acepta mensajes de tamaño arbitrario (hasta terabytes) y devuelve una etiqueta chica
(p.ej. 128 bits). El hash depende de \textbf{todos y cada uno} de los bits del mensaje
(efecto avalancha: cambiar 1-2 bits cambia totalmente la salida). Sirve como
identificador de contenido (deduplicación en backups, índice en redes P2P como
eDonkey/Kademlia, etc.). No permite reconstruir el mensaje original.

\paragraph{Colisiones.} Como hay más mensajes posibles que etiquetas, \textbf{las
colisiones existen forzosamente} (dos mensajes distintos con el mismo hash). Como no
hay clave, no se habla de ``falsificación'' sino de colisión.

\begin{definicion}{Tres propiedades de un hash criptográfico}
``Resistencia'' = la única forma de hallar lo buscado es por fuerza bruta (no hay
atajo algorítmico polinómico; es ``computacionalmente imposible'').
\begin{itemize}[leftmargin=1.4em,nosep]
  \item \textbf{Resistencia a preimagen:} dado un hash $y$, es computacionalmente
  imposible hallar un $x$ con $h(x)=y$.
  \item \textbf{Resistencia a segunda preimagen:} dado un mensaje $x$, es imposible
  hallar un segundo $x' \neq x$ con $h(x')=h(x)$.
  \item \textbf{Resistencia a colisión:} es imposible hallar \emph{dos} mensajes
  cualesquiera $x \neq x'$ con $h(x)=h(x')$. Es la más fuerte: implica las otras dos.
\end{itemize}
\end{definicion}

\paragraph{Usos.}
\begin{itemize}[leftmargin=1.4em,nosep]
  \item \emph{Resistencia a preimagen} $\to$ protocolos de \textbf{prueba de
  conocimiento}: demostrar que algo existe sin revelarlo (publicar el hash de una
  fórmula/patente y luego probar que se la tenía; juegos de cartas con compromiso).
  \item \emph{Resistencia a segunda preimagen} $\to$ índices de contenido /
  deduplicación: dos archivos con igual contenido tienen igual hash (match semántico).
\end{itemize}

\paragraph{Paradoja del cumpleaños.} Buscar una \textbf{colisión} es más fácil que una
preimagen: al ir generando mensajes se compara cada nuevo contra \emph{todos} los
anteriores, con complejidad del orden de la \textbf{raíz cuadrada} del espacio.
Medido en bits, el nivel de seguridad frente a colisiones es \textbf{la mitad}:
\begin{center}
\begin{tabular}{lll}
\toprule
Propiedad & Objetivo & Costo fuerza bruta \\
\midrule
Preimagen & dado $y$, hallar $x$: $h(x)=y$ & $\sim 2^{b}$ \\
2.\textsuperscript{a} preimagen & dado $x$, hallar $x' \neq x$ con igual hash & $\sim 2^{b}$ \\
Colisión & hallar cualquier par $x \neq x'$ colisionante & $\sim 2^{b/2}$ \\
\bottomrule
\end{tabular}
\end{center}
\noindent (con $b$ = bits de salida). Por eso el nivel de seguridad de una función se
mide por su resistencia a colisiones. Hoy el mínimo práctico de salida es 160 bits
(colisiones $> 2^{80}$).

\begin{center}
\begin{tikzpicture}[
  font=\footnotesize,
  arrow/.style={-{Stealth}},
  msg/.style={draw,rounded corners,fill=cyan!12,minimum width=10mm,minimum height=5mm}]
  % preimagen: comparo contra UN objetivo fijo
  \node[font=\bfseries] at (0,1.2) {Preimagen / 2.\textsuperscript{a} preimagen};
  \node[msg] (a1) at (-1,0.4) {$x_1$};
  \node[msg] (a2) at (-1,-0.3) {$x_2$};
  \node[msg] (a3) at (-1,-1.0) {$\vdots$};
  \node[draw,fill=yellow!25,rounded corners,minimum height=5mm] (y) at (1.6,-0.3) {$y$ fijo};
  \draw[arrow] (a1) -- (y);
  \draw[arrow] (a2) -- (y);
  \draw[arrow] (a3) -- (y);
  \node[align=center] at (0.3,-1.9) {comparo c/u contra\\ \emph{un} objetivo $\Rightarrow 2^{b}$};

  % colision: comparo cada nuevo contra TODOS los anteriores
  \node[font=\bfseries] at (6.6,1.2) {Colisión (cumpleaños)};
  \node[msg] (b1) at (5.2,0.4)  {$h(x_1)$};
  \node[msg] (b2) at (5.2,-0.3) {$h(x_2)$};
  \node[msg] (b3) at (5.2,-1.0) {$h(x_3)$};
  \node[msg] (b4) at (7.8,0.05) {$h(x_4)$};
  \draw[arrow] (b4) -- (b1);
  \draw[arrow] (b4) -- (b2);
  \draw[arrow] (b4) -- (b3);
  \node[align=center] at (6.6,-1.9) {cada nuevo contra \emph{todos}\\ los previos $\Rightarrow \sqrt{2^{b}}=2^{b/2}$};
\end{tikzpicture}
\end{center}
\diagcaption{Paradoja del cumpleaños: hallar una preimagen exige acertar contra un único objetivo ($\sim 2^{b}$), pero buscar una colisión compara cada hash nuevo contra todos los ya generados, por lo que basta del orden de la raíz cuadrada del espacio ($\sim 2^{b/2}$): el nivel de seguridad efectivo cae a la mitad.}

\paragraph{Modelo general iterativo (Merkle-Damgård, 1989).} La forma en que se
construyen la mayoría de las funciones (MD5, SHA-1, SHA-2): un \textbf{preprocesamiento}
ajusta el tamaño (padding + bloque con la longitud) y luego se aplica iterativamente una
\textbf{función de compresión} $f$ encadenando estados $H_i$ desde $H_0=IV$. La seguridad
de toda la función la da $f$.

\begin{center}
\begin{tikzpicture}[
  font=\footnotesize,
  arrow/.style={-{Stealth}},
  f/.style={draw,fill=blue!8,minimum height=9mm,minimum width=10mm}]
  \node (iv) at (-1.8,0) {$H_0=IV$};
  \node[f] (f1) at (0,0)   {$f$};
  \node[f] (f2) at (2.4,0) {$f$};
  \node[f] (f3) at (4.8,0) {$f$};
  \node     (d)  at (7.0,0) {$\cdots$};
  \node[f] (ft) at (8.6,0) {$f$};
  \node[draw,minimum height=7mm] (g) at (10.8,0) {$g$};
  \node[draw,fill=yellow!25,rounded corners] (out) at (10.8,-1.6) {$h(x)$};
  % encadenamiento horizontal
  \draw[arrow] (iv) -- (f1);
  \draw[arrow] (f1) -- (f2) node[midway,above]{$H_1$};
  \draw[arrow] (f2) -- (f3) node[midway,above]{$H_2$};
  \draw[arrow] (f3) -- (d);
  \draw[arrow] (d)  -- (ft);
  \draw[arrow] (ft) -- (g) node[midway,above]{$H_t$};
  \draw[arrow] (g)  -- (out);
  % bloques de mensaje entrando por arriba
  \node (x1) at (0,1.4)   {$x_1$};
  \node (x2) at (2.4,1.4) {$x_2$};
  \node (x3) at (4.8,1.4) {$x_3$};
  \node (xt) at (8.6,1.4) {$x_t$};
  \draw[arrow] (x1) -- (f1);
  \draw[arrow] (x2) -- (f2);
  \draw[arrow] (x3) -- (f3);
  \draw[arrow] (xt) -- (ft);
\end{tikzpicture}
\end{center}
\diagcaption{Merkle-Damgård: el mensaje formateado $x=x_1x_2\cdots x_t$ se procesa bloque a bloque por la función de compresión $f$, encadenando $H_i$ desde $H_0=IV$; una transformación final $g$ produce la salida $h(x)=g(H_t)$.}

\paragraph{Funciones en la práctica.}
\begin{itemize}[leftmargin=1.4em,nosep]
  \item \textbf{MD5}: \emph{quebrada}, no usar (aparece en protocolos viejos).
  \item \textbf{SHA-1} (160 bits): borderline, complicada.
  \item \textbf{SHA-2}: durante mucho tiempo la recomendada (3 variantes), pero $\sim$1 orden
  de magnitud más lenta que SHA-1.
  \item \textbf{SHA-3 (Keccak)}: ganador del concurso abierto, mismos tamaños que SHA-2
  pero velocidad de SHA-1 y estructura distinta. \textbf{Estándar recomendado hoy};
  256 bits ``alcanza y sobra''. SHA = Secure Hash Algorithm (estándares del NIST).
\end{itemize}

\begin{ojo}
Errores típicos de Verdadero/Falso (parciales 2025):
\begin{itemize}[leftmargin=1.4em,nosep]
  \item \textbf{``Un hash sirve para encriptar''} $\to$ \textbf{FALSO} (Recup. 2025). Un hash
  \emph{no es reversible}: no permite recuperar el mensaje, no encripta.
  \item \textbf{``MD5 es un criptosistema asimétrico''} $\to$ \textbf{FALSO} (2025-2). MD5 es
  una función de \emph{hash}, no un criptosistema asimétrico.
\end{itemize}
\end{ojo}

\subsection{HMAC: MAC a partir de una función de hash}

La segunda gran forma de construir un MAC infalsificable es a partir de una función de
hash libre de colisiones. Los HMAC suelen ser 1-2 órdenes de magnitud \textbf{más
rápidos} que un CBC-MAC (los hash están diseñados para procesar grandes volúmenes muy
rápido), por eso se prefieren a escala.

\begin{definicion}{HMAC (especificado en una RFC)}
Dada una función de hash $H$ y una clave $k$, con constantes públicas fijas
\texttt{ipad} y \texttt{opad}:
$$\mathrm{HMAC}_k(m) = H\big((k \oplus \texttt{opad}) \,\|\, H((k \oplus \texttt{ipad}) \,\|\, m)\big).$$
\end{definicion}

\begin{center}
\begin{tikzpicture}[
  font=\footnotesize,
  arrow/.style={-{Stealth}},
  xor/.style={draw,circle,inner sep=1pt},
  h/.style={draw,fill=blue!8,minimum height=8mm,minimum width=12mm}]
  % hash interno
  \node[xor] (xi) at (0,0) {$\oplus$};
  \node (ki) at (-1.6,0.4) {$k$};
  \node (ipad) at (-1.6,-0.4) {\texttt{ipad}};
  \draw[arrow] (ki) -- (xi);
  \draw[arrow] (ipad) -- (xi);
  \node (m) at (0,1.3) {$m$};
  \node[h] (hin) at (1.9,0) {$H$};
  \draw[arrow] (xi) -- (hin) node[midway,above]{$\|$};
  \draw[arrow] (m) -| (1.4,0.5) -- (hin);
  % hash externo
  \node[xor] (xo) at (4.4,0) {$\oplus$};
  \node (ko) at (4.4,1.0) {$k$};
  \node (opad) at (4.4,-1.0) {\texttt{opad}};
  \draw[arrow] (ko) -- (xo);
  \draw[arrow] (opad) -- (xo);
  \node[h] (hout) at (6.6,0) {$H$};
  \draw[arrow] (hin) -- (hout) node[midway,above]{interno};
  \draw[arrow] (xo) -- (hout) node[midway,above]{$\|$};
  \node[draw,fill=yellow!25,rounded corners] (t) at (9.2,0) {$\mathrm{HMAC}_k(m)$};
  \draw[arrow] (hout) -- (t);
\end{tikzpicture}
\end{center}
\diagcaption{HMAC: hash interno $H((k\oplus\texttt{ipad})\,\|\,m)$ y luego hash externo $H((k\oplus\texttt{opad})\,\|\,(\text{interno}))$. El mensaje se procesa una sola vez (el externo solo toca la etiqueta interna), por eso es muy eficiente.}

La clave concatenada como prefijo rompe la posibilidad de aprender estados
intermedios (misma idea que en CBC-MAC). El mensaje se procesa una sola vez (el hash
exterior solo procesa la etiqueta interna), por lo que es muy eficiente. Hay
\textbf{demostración formal}: si $H$ es libre de colisiones, HMAC es infalsificable.
En sistemas se ven HMAC-SHA1, HMAC-MD5 (evitar), HMAC-SHA3.

\begin{examen}
\textbf{Validar sin exponer el secreto (parcial 2017, caso CVV).} Se quiere verificar
un dato sensible (p.ej. el CVV de una tarjeta) sin almacenarlo ni transmitirlo en
claro: se usa un \textbf{HMAC con clave simétrica}. Se guarda/compara el tag en vez
del valor; quien no tiene la clave no puede generar ni validar el tag.
\end{examen}

\subsection{Cifrado autenticado: confidencialidad + integridad}

Cuando se necesitan \textbf{las dos cosas} (confidencialidad e integridad) hay tres
formas de combinar un criptosistema con un MAC. Dos son seguras, una no:

\begin{center}
\renewcommand{\arraystretch}{1.25}
\begin{tabular}{p{4.6cm}p{6.2cm}p{2.2cm}}
\toprule
\textbf{Construcción} & \textbf{Idea} & \textbf{¿Seguro?} \\
\midrule
\textbf{Encrypt-and-MAC} (en paralelo) & cifrar $m$ por un lado y etiquetar $m$ por el otro & \textbf{NO} \\
\textbf{MAC-then-Encrypt} (autenticar y luego cifrar) & etiquetar $m$, luego cifrar $m$ y el tag juntos & \emph{Puede} \\
\textbf{Encrypt-then-MAC} (cifrar y luego autenticar) & cifrar $m$, luego etiquetar el \emph{texto cifrado} $c$ & \textbf{SÍ} \\
\bottomrule
\end{tabular}
\end{center}

\begin{itemize}[leftmargin=1.4em]
  \item \textbf{Encrypt-and-MAC (paralelo) — inseguro:} no hay garantía de que la
  etiqueta del mensaje no revele información sobre $m$. La seguridad de un MAC cubre
  infalsificabilidad, \emph{no} confidencialidad. Publicar $t=\mathrm{Mac}(m)$ junto al
  cifrado puede filtrar información del plano.
  \item \textbf{MAC-then-Encrypt — puede ser seguro:} se cifra también el tag, así su
  posible filtración queda protegida. Lo usa SSH. Pero \textbf{no hay prueba genérica};
  solo pruebas para combinaciones concretas, y hay que tener cuidado si el MAC es un
  CBC-MAC con la misma primitiva de cifrado (simetrías explotables).
  \item \textbf{Encrypt-then-MAC — siempre seguro:} hay \textbf{demostración general}.
  Es lo que se llama \textbf{cifrado autenticado}.
\end{itemize}

\begin{center}
\begin{tikzpicture}[
  font=\footnotesize,
  arrow/.style={-{Stealth}},
  enc/.style={draw,fill=blue!8,minimum height=6mm,minimum width=11mm},
  mac/.style={draw,fill=green!10,minimum height=6mm,minimum width=11mm}]

  % ---- Encrypt-and-MAC (inseguro) ----
  \node[font=\bfseries] at (0,1.4) {Encrypt-and-MAC};
  \node[font=\bfseries\scriptsize,red] at (0,0.95) {INSEGURO};
  \node (m1) at (0,0.3) {$m$};
  \node[enc] (e1) at (-1.1,-0.8) {$\mathrm{Enc}_{k_1}$};
  \node[mac] (mc1) at (1.1,-0.8) {$\mathrm{Mac}_{k_2}$};
  \draw[arrow] (m1) -- (e1);
  \draw[arrow] (m1) -- (mc1);
  \node (c1) at (-1.1,-1.8) {$c$};
  \node (t1) at (1.1,-1.8) {$t$};
  \draw[arrow] (e1) -- (c1);
  \draw[arrow] (mc1) -- (t1);
  \node[align=center,font=\scriptsize] at (0,-2.5) {salida $(c,t)$\\ $t$ puede filtrar $m$};

  % ---- MAC-then-Encrypt (puede) ----
  \node[font=\bfseries] at (5,1.4) {MAC-then-Encrypt};
  \node[font=\bfseries\scriptsize] at (5,0.95) {puede ser seguro};
  \node (m2) at (5,0.3) {$m$};
  \node[mac] (mc2) at (5,-0.65) {$\mathrm{Mac}_{k_2}$};
  \draw[arrow] (m2) -- (mc2);
  \node[enc] (e2) at (5,-1.75) {$\mathrm{Enc}_{k_1}$};
  \draw[arrow] (mc2) -- (e2) node[midway,right]{$m\,\|\,t$};
  \node (c2) at (5,-2.7) {$c$};
  \draw[arrow] (e2) -- (c2);

  % ---- Encrypt-then-MAC (seguro) ----
  \node[font=\bfseries] at (10,1.4) {Encrypt-then-MAC};
  \node[font=\bfseries\scriptsize] at (10,0.95) {\textbf{SÍ} (cifrado autenticado)};
  \node (m3) at (10,0.3) {$m$};
  \node[enc] (e3) at (10,-0.65) {$\mathrm{Enc}_{k_1}$};
  \draw[arrow] (m3) -- (e3);
  \node[mac] (mc3) at (10,-1.75) {$\mathrm{Mac}_{k_2}$};
  \draw[arrow] (e3) -- (mc3) node[midway,right]{$c$};
  \node (t3) at (10,-2.7) {salida $(c,t)$};
  \draw[arrow] (mc3) -- (t3);
  \node[draw,rounded corners,fill=yellow!25,align=center,font=\scriptsize] at (10,-3.6)
    {MAC sobre el cifrado $c$\\ claves $k_1\neq k_2$};
\end{tikzpicture}
\end{center}
\diagcaption{Las tres construcciones. \textbf{Encrypt-and-MAC} (en paralelo): inseguro, el tag de $m$ puede filtrar información. \textbf{MAC-then-Encrypt}: se cifra $m\|t$, puede ser seguro pero sin prueba genérica. \textbf{Encrypt-then-MAC}: se etiqueta el \emph{texto cifrado} $c$ con claves independientes $k_1\neq k_2$; siempre seguro y es el ``cifrado autenticado''.}

\begin{definicion}{Cifrado autenticado (Encrypt-then-MAC)}
Requiere un criptosistema \textbf{CPA-Secure} y un MAC \textbf{infalsificable}, con
\textbf{claves independientes} ($k_1 \neq k_2$). Se construye un criptosistema nuevo:
\begin{itemize}[leftmargin=1.4em,nosep]
  \item $\mathrm{Gen}$: genera ambas claves ($k_1$ del cifrado, $k_2$ del MAC).
  \item Cifrado: se cifra $m$ obteniendo $c$, y se acompaña con la \textbf{etiqueta del
  texto cifrado} $t=\mathrm{Mac}_{k_2}(c)$.
  \item Descifrado/verificación: \textbf{primero se verifica} el tag sobre $c$; si pasa,
  se descifra; si no, la operación \textbf{falla} (error/excepción).
\end{itemize}
\end{definicion}

\begin{destacado}
\textbf{Encrypt-then-MAC = MAC sobre el CIFRADO $c$, con claves independientes
$k_1 \neq k_2$.} Aparece en parciales 2016, 2018, 2023. La novedad: si alguien
manipula el texto cifrado, el descifrado \textbf{falla} (antes solo quedaba ``basura''
según la maleabilidad).
\end{destacado}

\begin{examen}
\textbf{V/F (parciales 2018, 2023, 2016):} ``Cifrar y al mismo tiempo hacer MAC de $m$
con $k_1=k_2$'' (Encrypt-and-MAC con clave única) es \textbf{INSEGURO}
$\Rightarrow$ la afirmación de que es seguro es \textbf{FALSA}. El orden correcto es
MAC sobre el cifrado $c$, con claves \textbf{independientes}.
\end{examen}

\begin{examen}
\textbf{Construcción analizada a mano (parcial 2025-1): $c = E_{k_1}(m \,\|\, H(k_2\|m))$.}
Hay que describir los pasos del receptor: descifra $c$ con $k_1$, separa $m$ y el hash
adjunto, recomputa $H(k_2 \| m)$ con su clave $k_2$ y compara. Si coincide, el mensaje
es \textbf{íntegro y autenticado} (solo quien tiene $k_2$ pudo generar ese
``HMAC casero''). Da confidencialidad ($E_{k_1}$), integridad y autenticación, pero
\textbf{NO da no-repudio} (clave simétrica compartida).
\end{examen}

\paragraph{Modos AEAD en la práctica.} Modos de encadenamiento que convierten una
primitiva de bloque en cifrado autenticado, sin el overhead de doble función/doble
clave:
\begin{itemize}[leftmargin=1.4em,nosep]
  \item \textbf{CCM} (Counter with CBC-MAC): sigue la variante MAC-then-Encrypt (tiene
  prueba específica). Calcula el CBC-MAC y lo cifra junto al mensaje en modo Counter;
  pasa una sola vez por el mensaje pero procesa cada bloque dos veces. Usa una sola
  clave. Seguro si el IV y el nonce no se repiten.
  \item \textbf{GCM} (Galois Counter Mode): usa modo Counter (cifra cada bloque una
  vez) y combina los resultados con una función de hash rápida (GHASH, aritmética de
  polinomios). Mucho más rápido que CCM. \textbf{AES-GCM} es la recomendación práctica:
  rápido, estandarizado, libre de patentes, en todas las librerías serias.
\end{itemize}

\begin{destacado}
\prof{Si se tienen que llevar algo: en sistemas reales usen criptosistemas
autenticados.} En el 99\% de los casos, si se necesita confidencialidad también se
necesita integridad (el cifrado común puede ser manipulado). Recomendación concreta:
\textbf{AES en modo GCM}.
\end{destacado}

\subsection{No-repudio: por qué los MAC no lo dan}

\begin{ojo}
\textbf{ERROR TÍPICO MUY IMPORTANTE (parciales 2025-1 Ej.3d, 2025-2 Ej.5b):}
afirmar que un \textbf{MAC/HMAC da NO-REPUDIO}.\\[4pt]
\textbf{NO lo da.} La clave simétrica es \textbf{compartida}: ambas partes (emisor y
receptor) tienen la misma clave y \emph{ambas} pueden generar un tag válido. Por lo
tanto, el receptor no puede probar ante un tercero que el mensaje lo originó el
emisor (él mismo pudo haberlo generado). El MAC da \textbf{integridad} y
\textbf{autenticación de origen} entre las partes que comparten la clave, pero
\textbf{no} no-repudio.\\[4pt]
El \textbf{no-repudio requiere firma digital} (criptografía \textbf{asimétrica}:
clave privada que solo posee el firmante).
\end{ojo}

\end{document}
